% Fragmento público generado desde propietarios íntegros.
% Residencia: 52.5.2.
% Fuente: ../incoming/radion_monografia_20260827/manuscrito/sections/06_interior_helice.tex:109-201
\subsubsection{Helical elevation of the nonadic return}

\paragraph{Nonadic decomposition and memory quotient}

The elementary operation
\begin{equation}
n=9q_9(n)+\rho_9(n),
\qquad
H_9(n)=(\rho_9(n),q_9(n))
\label{eq:H9}
\end{equation}
satisfies
\begin{equation}
H_9(n+9)=H_9(n)+(0,1).
\label{eq:H9-helix}
\end{equation}
The first coordinate closes; the second ascends. This identity is the
minimal arithmetic version of a screw.

For a spectral fiber,
\[
n\longmapsto(\tau_\gamma^n,n)
\]
is a helix over \(S^1\). Adding the contraction \((5/9)^n\) produces an
inner spiral. The radion participates in both readings: a unit of phase
on the boundary and an oriented unit of action in the lift.

\paragraph{Memory range as dimensional reader}

Suppose that a visible projection preserves the phase value but loses \(d\) linearly independent memory cocycles. Any faithful lift requires at least \(d\) additional coordinates to reconstruct them. This motivates a precise reading of Dimensional Mechanics:
\begin{equation}
\begin{aligned}
\text{emergent dimension}
&=\text{minimum rank of independent memory}\\
&\phantom{=}\text{not publishable on the visible sheet}.
\end{aligned}
\label{eq:dimension-memory}
\end{equation}
It is not asserted that every physical dimension is a counter. Rather, in the
TPK domain, every independent memory that must survive forces a
coordinate of the lift.

\begin{narrativebox}
Dimension appears when the projection can no longer preserve the entire
history. The circle suffices to state where the phase is; it does not
suffice to state how many times it returned, which sheet it transported,
which interval contracted, or which carry survived.
\end{narrativebox}

\paragraph{Cylinder Assembly, Spirals, and Helical Trajectories}

The coinductive cylinders are not physical tubes added to a spiral. They are
nested intervals that preserve prefixes and narrow the value. The spiral
describes transverse contraction; the helix describes longitudinal memory;
the cylinder describes the set of compatible values at a finite
depth. In the limit, the intersection of the cylinders selects a unique
section.

The spatial image brings together, without confusion:
\[
\begin{array}{ccl}
\text{angle} &\leftrightarrow& \text{visible phase},\\
\text{height} &\leftrightarrow& \text{return/memory},\\
\text{cylinder radius} &\leftrightarrow& \text{refinement uncertainty},\\
\text{sheet} &\leftrightarrow& \text{bidirectional orientation}.
\end{array}
\]
The rate \(3^{-54}\) contracts the cylinder by return; it is not added to the angle and
does not replace the actual carry \(\epsR\).

\subsubsection{Phase closures of orders \(2\pi\) and \(4\pi\)}

In the circular projection, \(2\pi\) restores the phase. In an oriented lift
with two sheets, one turn may close on the conjugate sheet and require a
second turn to restore the orientation:
\[
U_{2\pi}=-I,\qquad U_{4\pi}=I.
\]
The structure agrees with spinorial periodicity when the
functor is applied to a representation \(SU(2)\); the TPK lift is not identified with
\(SU(2)\) without that map.

The areal reading is consistent:
\[
\mathcal A(2\pi)=\hfull,\qquad
\mathcal A(4\pi)=2\hfull.
\]
In a Möbius band, one visible circuit reverses the sheet and two circuits
restore the orientation. The \(2\pi/4\pi\) is not an accidental duplication;
it is the difference between phase return and return of the enriched state.
